\documentclass[12pt]{amsart}
\pagestyle{empty}
\usepackage[hmargin=1in,vmargin=1in]{geometry}
\usepackage{url,  mathrsfs}
\usepackage{hyperref}
\usepackage{fancyhdr} 
\usepackage[T1]{fontenc}


\usepackage{url,  mathrsfs}
\usepackage{hyperref}
\usepackage{amsmath}
\usepackage{graphicx,  epstopdf}
\usepackage{subfig}
\usepackage{color}
\usepackage{bbm}
\usepackage{subfloat}
\usepackage[draft]{fixme}
\usepackage{bm} 
\usepackage{bbm} 
\usepackage{epigraph}
\usepackage{endnotes}
\usepackage{ifpdf}
\usepackage{array}

\usepackage{multicol}
\usepackage{multirow}
\usepackage{graphicx}
\usepackage{tikz}

\newcommand{\A}{\mathbb{A}}
\newcommand{\Q}{\mathbb{Q}}
\newcommand{\N}{\mathbb{N}}
\newcommand{\Z}{\mathbb{Z}}
\newcommand{\R}{\mathbb{R}}
\newcommand{\F}{\mathbb{F}}
\newcommand{\C}{\mathbb{C}}
\renewcommand{\P}{\mathbb{P}}
\newcommand{\I}{\mathcal{I}}
\newcommand{\V}{\mathcal{V}}
\newcommand{\cH}{\mathcal{H}}
\newcommand{\ol}{\overline}
\newcommand{\J}{\mathcal{J}}
\newcommand{\X}{\mathcal{X}}
\theoremstyle{definition}
\newtheorem{prob}{Problem}
\newtheorem{extraProb}{Extra Problem}

\DeclareMathOperator{\conv}{conv}
\DeclareMathOperator{\cone}{cone}
\DeclareMathOperator{\ext}{ext}




\begin{document}
 \thispagestyle{empty}
\begin{center} {\large \textbf{Math 380 -- Homework 8}} \smallskip \\ 
Due on Thursday, June 4 \\ \ \\
 \end{center}

You are welcome to talk with other students in the class about problems but should write up solutions on your own. 
Solutions can be handwritten or typed but need to be legible and submitted via Gradescope by the end of the day 
on Thursday. You should justify all your answers in order to receive full credit.  
\bigskip 

\begin{flushleft}

{\bf Good practice problems} (not to be turned in): \\

CLO Section 4.5, Exercises 2, 3, 4, 7, 8, 9 \\
CLO Section 4.6, Exercises 3, 4, 8\\
\ \ \\
\bigskip 




\begin{prob} CLO \S4.5, Exercise 11  \smallskip  \\
If $f \in \mathbb C[x_1,\dots,x_n]$ is irreducible, then $\mathbf V(f)$ is irreducible.
\end{prob}

Consider the ideal $I = \langle f \rangle \subseteq \mathbb C [x_1,\dots,x_n]$. We can show this ideal is prime. Let $g,h \in \mathbb C[x_1,\dots,x_n]$, and suppose that $g \cdot h \in I$. This means that $f$ divides $g \cdot h$, and because $\mathbb C[x_1,\dots,x_n]$ is a UFD and $f$ is irreducible, either $f \mid g$ or $f \mid h$. Thus, either $g \in I$ or $h \in I$, completing the proof that $I$ is prime.
\\ \ \\
Since every prime ideal is radical (see the problem below), $I = \sqrt{I} = \mathbf I(\mathbf V(I))$ by the Nullstellensatz as $\mathbb C$ is algebraically closed. Considering the variety $V = \mathbf V(I)$, we have that $\mathbf I(V)$ is prime, so $V=\mathbf V(f)$ must be irreducible as desired.

\bigskip 

\begin{prob}  CLO \S4.6, Exercise 1 \smallskip \\
The intersection of any collection of prime ideals is radical.
\end{prob}
First, we can show that every prime ideal is radical. Let $P \subseteq k[x_1,\dots,x_n]$ be any prime ideal, and let $f \in k[x_1,\dots,x_n]$. We can show by induction that $f^n \in P \implies f \in P$ for all $n \in \mathbb Z_{\ge 1}$. The base case of $n=1$ is a tautology. For the inductive case, suppose that we had $f^{n+1} = f\cdot f^n \in P$. Since $P$ is a prime ideal, either $f \in P$ or $f^n \in P$. If $f^n \in P$, by the inductive hypothesis that $f^n \in P \implies f\in P$, we must have that $f \in P$. Thus, in either case, we have $f \in P$, and we are done. $P$ is a radical ideal.
\\ \ \\
Next, we can show that the intersection of any family of radical ideals is radical. Let $I_j$ be a family of radical ideals index by $j \in \mathcal J$, and let $I = \bigcap_{j \in \mathcal J} I_j$. Let $g \in k[x_1,\dots,x_n]$ be any polynomial. Clearly, $I$ is an ideal since \[g \in I \iff \forall j \in \mathcal J, g \in I_j\]
and every $I_j$ is an ideal. Suppose that $g^n \in I$, or $\forall j \in \mathcal J, g^n \in I_j$. Since every $I_j$ is a radical ideal, $\forall j \in \mathcal J, g \in I_j$ so $g \in I$. Thus, $I$ is a radical ideal.

\newpage{}
\begin{prob}  CLO \S4.6, Exercise 10 \smallskip \\
Let $f \in \mathbb C[x_1,\dots,x_n]$, and let $f = f_1^{a_1}\dots f_r^{a_r}$ be a decomposition of $f$ into distinct irreducible factors. Then $\mathbf V(f)=\mathbf V(f_1)\cup \dots \cup \mathbf V(f_r)$ is a decomposition of $\mathbf V(f)$ into irreducible varieties and $\mathbf I(\mathbf V(f)) = \langle f_1\dots f_r\rangle$.   
\end{prob}
\bigskip 
First, note that $\langle f\rangle \subseteq \langle f_i\rangle$ for all $1 \le i \le r$ as every $f_i$ divides $f$. Thus, $\mathbf V(f) \supseteq \mathbf V(f_i)$ for every $i$, and so $\mathbf V(f) \supseteq \mathbf V(f_1)\cup \dots\cup \mathbf V(f_r)$. We can also show that every point in $\mathbf V(f)$ is contained in some $\mathbf V(f_i)$. Let $p \in \mathbf V(f)$, and note that $f(p)=f_1^{a_1}(p)\dots f_r^{a_r}(p) = 0$. Because $\mathbb C$ is an integral domain, we must have $f_i^{a_i}(p) = 0$ for some $i$, and in fact $f_i(p)=0$. Thus, $p \in \mathbf V(f_i)$ for some $i$. We now have that $\mathbf V(f) \subseteq \mathbf V(f_1)\cup \dots \cup \mathbf V(f_r)$, and $\mathbf V(f) = \mathbf V(f_1)\cup \dots \cup \mathbf V(f_r)$. Because every $f_i$ is an irreducible polynomial in $\mathbb C[x_1,\dots,x_n]$, by the previous problem we have that $\mathbf V(f_i)$ must be an irreducible variety. This is a decomposition of $\mathbf V(f)$ into irreducible varieties.
\\ \ \\
We are working over an algebraically closed field, so by the Nullstellensatz we have $\mathbf I(\mathbf V(f)) = \sqrt{\langle f\rangle}$. Let $f_\text{red} = f_1\dots f_r$, and note that $\langle f_\text{red}\rangle \subseteq \sqrt{\langle f\rangle}$ as $f_\text{red}^{\max\{a_1,\dots,a_r\}} \in \langle f\rangle$. We can also show that $\sqrt{\langle f \rangle}\subseteq \langle f_\text{red}\rangle$. Let $g\in \sqrt{\langle f \rangle}$, and let $n \in \mathbb Z_{\ge 1}$ be such that $g^n \in \langle f\rangle$. We can write 
\[g^n = f_1^{a_1}\dots f_r^{a_r} \cdot h\]
for some $h \in \mathbb C[x_1,\dots, x_n]$. Thus, every $f_i$ is an irreducible polynomial that divides $g^n$, and so each $f_i$ divides $g$. Because these are all distinct irreducible elements in a UFD, $f_\text{red} = f_1\dots f_r$ must also divide $g$, and $g \in \langle f_\text{red}\rangle$. Thus, $\sqrt{\langle f \rangle} \subseteq \langle f_\text{red}\rangle$, and $\mathbf I(\mathbf V(f)) = \langle f_1\dots f_r\rangle$.
\end{flushleft}
\end{document}